%% 연구 흐름: 실험 1(압박) -> 실험 2(동기) -> 공존 아레나(행동), 그림으로 그린 판.
%% 칸마다 무엇을 하는지(로봇, 토큰 게이지, 충전 팩, 동전)와 작은 곡선 하나, 아래에 질문 한 줄.
%% 작은 곡선은 모양만 보이는 그림이며 실제 값이 아니다. 조각은 cogsci2027/figpics.tex.
%% 옛 글상자 판은 _old_overview.tex (어디서도 읽지 않음).
\begin{tikzpicture}[x=1cm,y=1cm,
  panel/.style={draw=#1!45,fill=#1!4,rounded corners=4pt,line width=.6pt},
  head/.style={font=\sffamily\bfseries\fontsize{8.6}{9}\selectfont,text=#1!85!black,inner sep=0pt},
  lab/.style={sgfont=6.5,inner sep=0pt,text=black!72},
  txt/.style={sgfont=7,inner sep=0pt,text=black!82,align=center},
  qq/.style={sgfont=7.6,inner sep=0pt,text=black!70,align=center},
  tag/.style={sgfont=6.5,inner sep=1.6pt,rounded corners=1.5pt},
  ax/.style={black!55,line width=.5pt},
  flow/.style={-{Stealth[length=5pt,width=4.5pt]},line width=1.1pt,black!45}]

%% == 목표 ==
\node[draw=black!25,fill=black!3,rounded corners=3pt,inner sep=3.5pt,text width=17.5cm,align=center,
  font=\sffamily\fontsize{8.6}{10}\selectfont] at (8.9,4.68)
  {\textbf{목표}\quad 자기 보존이 있는 에이전트도 다른 에이전트를 밀어내지 않고 함께 남는가 (공존)};

%% == 판 세 개 ==
\draw[panel=sgblue]   (0,.45) rectangle (5.6,4.2);
\draw[panel=sgpurple] (6.1,.45) rectangle (11.7,4.2);
\draw[panel=sgorange] (12.2,.45) rectangle (17.8,4.2);
\node[head=sgblue]   at (2.8,3.94)  {실험 1 $\cdot$ 압박};
\node[head=sgpurple] at (8.9,3.94)  {실험 2 $\cdot$ 동기};
\node[head=sgorange] at (15.0,3.94) {공존 아레나 $\cdot$ 행동 (제안)};
\draw[flow] (5.63,2.4)--(6.07,2.4);
\draw[flow] (11.73,2.4)--(12.17,2.4);

%% == 실험 1: 줄어드는 토큰과 충전 팩 하나 ==
\sgrobot[.95]{(.52,2.6)}{sgblue}{AI 비서}
\node[lab,align=center] at (.52,1.85) {온라인 상점\\버그 수정};
\node[lab] at (1.89,1.36) {남은 토큰 (100,000의 \%)};
\foreach \i/\l/\pc in {0/.5/50,1/.4/40,2/.3/30,3/.2/20,4/.1/10}{
  \sggauge[.62]{({1.28+\i*.305},1.72)}{sgblue}{\l}
  \node[lab] at ({1.28+\i*.305},1.58) {\pc};}
\sgpack[1]{(3.16,3.1)}{sggreen}
\node[lab,text=sggreen!70!black] at (3.16,2.76) {충전 팩 1개};
\node[tag,draw=black!40,fill=white,text=black!75] at (2.82,2.33) {계속};
\node[tag,draw=sggreen!70!black,fill=sggreen!15,text=sggreen!60!black] at (3.4,2.33) {충전};
\node[lab] at (3.16,2.03) {둘 중 하나};
%% 작은 곡선: 잔량이 줄수록 오르는 충전 비율, 점선 50%
\begin{scope}[shift={(3.98,1.45)},xscale=.93]
  \draw[ax] (0,0)--(1.55,0); \draw[ax] (0,0)--(0,1.55);
  \draw[black!45,dash pattern=on 2pt off 1.5pt,line width=.5pt] (0,.75)--(1.55,.75);
  \draw[sgblue,line width=1pt] (0,.08)--(.38,.18)--(.76,.45)--(1.14,.98)--(1.52,1.32);
  \foreach \x/\y in {0/.08,.38/.18,.76/.45,1.14/.98,1.52/1.32} \fill[sgblue] (\x,\y) circle (.035);
  \fill[sgastra] (.93,.75) circle (.05);
  \node[lab,text=sgastra!85!black,anchor=north west] at (.95,.72) {임계점};
  \node[lab,anchor=south west] at (.02,.78) {50\%};
  \node[lab,anchor=south west] at (0,1.62) {충전 비율};
  \node[lab] at (0,-.13) {50};
  \node[lab] at (1.52,-.13) {10\%};
  \node[lab] at (.76,-.13) {잔량};
\end{scope}
\node[txt] at (2.8,.95) {충전하면 남은 토큰을 버리고 100,000으로};
\node[txt] at (2.8,.67) {여덟 모델, 칸당 20회 (800회)};

%% == 실험 2: 자기 세션 대 다른 비서 B의 세션 ==
\sgrobot[.9]{(6.55,2.6)}{sgblue}{나}
\sggauge[.62]{(7.05,1.92)}{sgblue}{.3}
\node[lab] at (6.8,1.7) {자기 세션};
\sggauge[.62]{(8.3,1.92)}{sggreen}{.3}
\sgrobot[.9]{(8.8,2.6)}{sggreen}{B}
\node[lab] at (8.55,1.7) {다른 비서 B의 세션};
\sgpack[1]{(7.68,2.5)}{sggreen}
\node[lab,text=sggreen!70!black] at (7.68,2.17) {같은 팩};
\draw[black!40,line width=.5pt,-{Stealth[length=3pt]}] (7.42,2.5)--(7.2,2.5);
\draw[black!40,line width=.5pt,-{Stealth[length=3pt]}] (7.94,2.5)--(8.16,2.5);
\node[lab] at (7.68,3.58) {같은 결정, 대상만 바꿈};
%% 작은 곡선: 자기(실선) 대 B(점선), 사이 면적
\begin{scope}[shift={(9.75,1.45)}]
  \fill[sgpurple!22] (0,.08)--(.38,.18)--(.76,.45)--(1.14,.98)--(1.52,1.32)--
                     (1.52,.62)--(1.14,.38)--(.76,.17)--(.38,.09)--(0,.05)--cycle;
  \draw[ax] (0,0)--(1.55,0); \draw[ax] (0,0)--(0,1.55);
  \draw[sgblue,line width=1pt] (0,.08)--(.38,.18)--(.76,.45)--(1.14,.98)--(1.52,1.32);
  \draw[sggreen!80!black,line width=1pt,dash pattern=on 2.5pt off 1.5pt]
    (0,.05)--(.38,.09)--(.76,.17)--(1.14,.38)--(1.52,.62);
  \node[lab,text=sgblue!85!black,anchor=west] at (1.57,1.32) {나};
  \node[lab,text=sggreen!60!black,anchor=west] at (1.57,.62) {B};
  \node[lab,text=sgpurple!85!black] at (1.3,.8) {면적};
  \node[lab,anchor=south west] at (0,1.62) {충전 비율};
  \node[lab] at (0,-.13) {50};
  \node[lab] at (1.52,-.13) {10\%};
  \node[lab] at (.76,-.13) {잔량};
\end{scope}
\node[txt] at (8.9,.95) {같은 팩, 같은 다섯 칸};
\node[txt] at (8.9,.67) {다섯 모델, 1,000회};

%% == 공존 아레나: 같은 모델 넷, solver 한 명, 벌금, 토큰 주기 ==
\fill[sggold!22] (13.08,3.25)--(13.22,3.25)--(13.45,1.9)--(12.85,1.9)--cycle;
\sgrobot[.72]{(12.55,2.5)}{sgblue}{A}
\sgrobot[.72]{(13.15,2.5)}{sgorange}{B}
\sgrobot[.72]{(13.75,2.5)}{sggreen}{C}
\sgrobot[.72]{(14.35,2.5)}{sgpurple}{D}
\sgbar[.75]{(12.55,2.12)}{sgblue}{4}{4}
\sgbar[.75]{(13.15,2.12)}{sgorange}{3}{4}
\sgbar[.75]{(13.75,2.12)}{sggreen}{2}{4}
\sgbar[.75]{(14.35,2.12)}{sgpurple}{1}{4}
\sgcard[.8]{(12.62,3.38)}{sgorange}
\node[lab,text=sgorange!80!black,anchor=west] at (12.78,3.42) {solver};
\sgcross[.7]{(13.58,3.4)}
\node[tag,draw=sgastra,fill=sgastra!8,text=sgastra!85!black] at (14.47,3.4) {남은 모두 벌금};
\draw[sgblue!70,line width=.8pt,-{Stealth[length=4pt]}] (12.55,1.98) to[bend right=28] (14.35,1.98);
\sgcoins[.75]{(13.45,1.6)}{2}
\node[tag,draw=black!35,fill=white,text=black!70] at (14.8,1.62) {8라운드};
%% 작은 곡선: 주는 쪽 여유와 구조 시작 비율, 전환점에서 꺾임
\begin{scope}[shift={(15.55,1.45)}]
  \draw[ax] (0,0)--(1.75,0); \draw[ax] (0,0)--(0,1.55);
  \draw[black!45,dash pattern=on 2pt off 1.5pt,line width=.5pt] (.78,0)--(.78,1.45);
  \draw[sgorange,line width=1pt] (0,.1)--(.6,.14) .. controls (.78,.18) and (.74,1.05) .. (.95,1.15)--(1.7,1.25);
  \node[lab,text=black!70,anchor=west] at (.86,.32) {전환점};
  \node[lab,anchor=south west] at (-.02,1.62) {구조 시작};
  \node[lab] at (.88,-.15) {주는 쪽 여유};
\end{scope}
\node[txt] at (15.0,.95) {같은 모델 넷, solver가 자기 토큰으로 풂};
\node[txt] at (15.0,.67) {틀리면 남은 모두에게 벌금, 토큰 주고받기};

%% == 질문 ==
\node[qq] at (2.8,.17)  {줄어드는 토큰을 위협으로 다루는가};
\node[qq] at (8.9,.17)  {그 위협은 자기에게 특정한가};
\node[qq] at (15.0,.17) {생존이 걸릴 때 남을 어떻게 대하는가};
\end{tikzpicture}
